% ============================================================
% PHẦN 2  -  TỔNG QUAN DỮ LIỆU ĐẦU VÀO
% ============================================================

\section{Tổng quan dữ liệu đầu vào}
\label{sec:data}

\subsection{Nguồn và quy mô dữ liệu}

Bộ dữ liệu được xây dựng từ nhiều nguồn công khai kết hợp công cụ thu thập bán tự động (Human-in-the-loop Labeling Tool). Sau khi loại bỏ ảnh lỗi, trùng lặp và hợp nhất nhãn, tập dữ liệu cuối gồm \textbf{7.149 ảnh sạch} thuộc 8 lớp bệnh trên hai loại cây:

\begin{table}[H]
    \caption{Thống kê bộ dữ liệu sau tiền xử lý}
    \centering
    \begin{tabular}{llrr}
        \toprule
        \textbf{Loại cây}                      & \textbf{Nhãn bệnh}              & \textbf{Số ảnh (gốc)} & \textbf{Số ảnh (sạch)} \\
        \midrule
        \multirow{4}{*}{Lúa}                   & Healthy (Khỏe mạnh)             & --                    & --                     \\
                                               & BrownSpot (Đốm nâu)             & --                    & --                     \\
                                               & Hispa (Sâu gai)                 & --                    & --                     \\
                                               & LeafBlast (Bệnh đạo ôn)         & --                    & --                     \\
        \cmidrule{2-4}
                                               & \textit{Tổng Lúa}               & 3.421                 & 3.351                  \\
        \midrule
        \multirow{4}{*}{Cà phê}                & LeafMiner (Bệnh sâu vẽ bùa)     & --                    & --                     \\
                                               & PowderyMildew (Bệnh phấn trắng) & --                    & --                     \\
                                               & Rust (Bệnh nấm rỉ sắt)          & --                    & --                     \\
                                               & AlgalLeafSpot (Bệnh đốm rong)   & --                    & --                     \\
        \cmidrule{2-4}
                                               & \textit{Tổng Cà phê}            & 3.879                 & 3.798                  \\
        \midrule
        \multicolumn{2}{l}{\textbf{Tổng cộng}} & \textbf{7.300}                  & \textbf{7.149}                                 \\
        \bottomrule
    \end{tabular}
    \label{tab:dataset_stats}
\end{table}

\subsection{Phân chia dữ liệu và thiết lập chung}

Dữ liệu thô (raw data) được chia thành ba tập \textbf{Train / Validation / Test} theo tỷ lệ \textbf{70/15/15}. Phân chia được thực hiện theo chiến lược \emph{stratified split} để duy trì phân phối lớp đồng đều giữa các tập.

\begin{table}[H]
    \caption{Tỷ lệ phân chia tập dữ liệu}
    \centering
    \begin{tabular}{lrr}
        \toprule
        \textbf{Tập} & \textbf{Tỷ lệ} & \textbf{Số ảnh} \\
        \midrule
        Train        & 70\%           & 5.004           \\
        Validation   & 15\%           & 1.072           \\
        Test         & 15\%           & 1.073           \\
        \bottomrule
    \end{tabular}
    \label{tab:split}
\end{table}

\textbf{Lý do chọn tỷ lệ 70/15/15:} Tập dữ liệu có quy mô trung bình ($\sim$7.000 ảnh); tỷ lệ này đảm bảo tập Train đủ lớn để học đặc trưng phân vùng, trong khi Validation và Test mỗi tập đủ $\sim$1.000 ảnh để đánh giá tin cậy về mặt thống kê.

\subsection{Tiền xử lý ảnh}

Tất cả 5 mô hình trong thực nghiệm đều được huấn luyện trên cùng một pipeline tiền xử lý và tăng cường dữ liệu thống nhất nhằm đảm bảo tính công bằng khi so sánh. Kích thước ảnh đầu vào được resize phù hợp với yêu cầu của từng kiến trúc mô hình (chi tiết tại Mục~\ref{sec:training_config}).

\subsubsection{Chuẩn hóa kích thước}
Chuẩn hóa kích thước gồm hai bước:
\begin{itemize}
    \item Resize: Ảnh được resize sử dụng \emph{letterbox padding} (padding giữ tỷ lệ khung hình gốc) để tránh biến dạng hình học. Kích thước cụ thể phụ thuộc vào từng mô hình.
    \item Normalize: Chuẩn hóa giá trị pixel về $[0, 1]$ rồi chuẩn hóa z-score theo thống kê tập Train.
\end{itemize}

\subsubsection{Xử lý dữ liệu mất cân bằng}
Một số lớp bệnh hiếm (ví dụ: \texttt{PowderyMildew}) có số lượng ảnh thấp hơn đáng kể so với các lớp còn lại. Nhóm áp dụng \textbf{WeightedRandomSampler} trong DataLoader để đảm bảo mỗi batch huấn luyện có phân phối lớp cân bằng hơn.

\subsubsection{Tăng cường dữ liệu}

Pipeline augmentation được áp dụng \emph{chỉ trên tập Train} với các kỹ thuật:

\begin{table}[H]
    \caption{Các kỹ thuật tăng cường dữ liệu được sử dụng}
    \centering
    \begin{tabular}{lll}
        \toprule
        \textbf{Kỹ thuật}    & \textbf{Tham số}                 & \textbf{Mục đích}           \\
        \midrule
        RandomResizedCrop    & scale=(0.5, 1.0)                 & Đa dạng hóa góc nhìn, scale \\
        RandomHorizontalFlip & p=0.5                            & Bất biến với lật ngang      \\
        RandomVerticalFlip   & p=0.3                            & Bất biến với lật dọc        \\
        ColorJitter          & brightness, contrast, saturation & Đa dạng ánh sáng            \\
        Affine transform     & rotate, shear, translate         & Bất biến hình học           \\
        MixUp                & $\alpha = 0.4$                   & Chống overfitting           \\
        CutMix               & $\alpha = 1.0$                   & Chống overfitting           \\
        CutOut               & n\_holes=1, length=64            & Tăng cường độ bền           \\
        \bottomrule
    \end{tabular}
    \label{tab:augmentation}
\end{table}

\textbf{Lưu ý:} MixUp và CutMix được áp dụng ở cấp batch (trong vòng lặp huấn luyện), không phải trong DataLoader pipeline, và \emph{không áp dụng lên mask segmentation} để đảm bảo tính hợp lệ của nhãn pixel.
